\section{Access Control}\label{sec:appendix:access_control}
\vspace{7.5cm}

Samarati and De Capitani di Vimercati~\cite{samarati_access_2000} defined \gls{ac} as ``the process of mediating every request to resources maintained by a system and determining whether the request should be granted or denied''. A resource usually consists of data such as messages or files. \gls{ac} is traditionally divided into three levels:
\begin{itemize}
	\item\textbf{Policy}: this abstract level consists of the rules stating which users can perform which operations on which resources. The policy is usually defined by the owner of the resources or the system (e.g., the administrator);
	\item\textbf{Model}: this intermediate level is a formal representation of the policy (e.g., \gls{rbac}~\cite{sandhu_access_1998} and \gls{abac}~\cite{horwitz_toward_2002} are two popular models) giving the semantics for granting or denying users' requests;
	\item\textbf{Enforcement}: this concrete level corresponds to the hardware and software entities that enforce the policy based on the chosen model.
\end{itemize}

Below, we briefly describe \gls{rbac} (see \Cref{sec:appendix:access_control:rbac}) and \gls{abac} (see \Cref{sec:appendix:access_control:abac}).


\subsection{Role-based Access Control}\label{sec:appendix:access_control:rbac}

\sbnote{TODO}


\subsection{Attribute-based Access Control}\label{sec:appendix:access_control:abac}

\sbnote{TO IMPROVE}

\gls{abac} is a fine-grained \gls{ac} model which regulates access to resources based on the evaluation of ``attributes'' relevant to making \gls{ac} decisions \cite{horwitz_toward_2002}. In detail, \gls{abac} may include three types of attributes: 
\begin{itemize}
	\item \textbf{subject attributes} define the characteristics of users (e.g., name, organization, role);
	\item \textbf{object attributes} define the characteristics of resources (e.g., title, size, category);
	\item \textbf{environmental attributes} describe the context in which a user's request occurs (e.g., date, network security level, timezone).
\end{itemize}

Attributes are also used as building blocks for \textbf{access structures}. Informally, an access structure implements a (rule of a) policy by specifying which subject, object and environmental attributes are required to execute an operation over a specific resource. More formally, given a set of attributes $\Attributes$, an access structure $\abeAccessStructure$ is a subset of the power set of $\Attributes$ minus the empty set, i.e., $\abeAccessStructure \subseteq (\mathcal{P}(\Attributes) \setminus \emptyset)$. Each element in $\abeAccessStructure$ is called \textbf{authorized attribute set} and it is said to \textbf{satisfy} $\abeAccessStructure$, whereas $(\mathcal{P}(\Attributes) \setminus \abeAccessStructure)$ identifies the \textbf{unauthorized} attribute sets. 

For brevity and ease of use, access structures are usually compiled and expressed as boolean formulas over attributes \cite{venema_systematizing_2023}.\footnote{Such boolean formulas may also be seen as \textit{policy trees} in which leaves represent attributes while internal nodes represent boolean operators (e.g., {\tt AND}, {\tt OR}, and {\tt NOT}).} For instance, consider a policy stating that the ECG data of a patient can be accessed only by either cardiologists or surgeons. Given a set of attributes $\Attributes = \{\mathit{cardiologist}, \mathit{nurse}, \mathit{surgeon}\}$, the access structure $\abeAccessStructure$ implementing the policy can be expressed either as the set of authorized attribute sets $\{ \{\mathit{cardiologist}\}, \{\mathit{surgeon}\}, \{\mathit{nurse}, \mathit{surgeon}\}, ... \}$ or, more concisely, as the boolean formula ``$\mathit{cardiologist} \lor \mathit{surgeon}$''. Compiling access structures into formulas also allows expressing conditions over attributes and their values (e.g., using $<, \leq, \ge, >, =$ and $\neq$).

%For the purposes of our work --- as we discuss at the end in \Cref{sec:background:abe} --- we consider subject attributes only hereafter. Hence, the state of an \gls{abac} policy $\mathbf{P}$ can be described as a tuple $\langle \Users, \Attributes, \Resources, \UA, \FA \rangle$, where $\Users$ is the set of users, $\Attributes$ is the set of (subject) attributes and $\Resources$ is the set of resources. $\UA \subseteq \Users \times \Attributes$ is the set of user-attribute assignments, where each element (or relation) in $\UA$ may be enriched with a value --- represented as an arbitrary binary string $\{0,1\}^*$ --- for the corresponding attribute; for instance, a user $u$ can be assigned to an attribute $a = ``\mathit{age}$'' with value $\attributeValue = ``\mathit{10010}$'' (i.e., 18). Then, $\FA \subseteq \Permissions \times (\mathcal{P}(\Attributes) \setminus \emptyset)$ is the set of permission-access structure assignments, being $\Permissions \subseteq \Resources \times \Operations$ a derivative set of $\Resources$ combined with a fixed set of operations $\Operations$ (e.g., read, write). Both $\Operations$ and $\Permissions$ are not part of the state of the \gls{ac} policy, as $\Operations$ remains constant over time and $\Permissions$ is derivative of $\Resources$ and $\Operations$. A user $u$ can use a permission $\langle f, \mathit{op} \rangle$ if $\exists (f, \mathit{op}, \Attributes') \in \FA \sth \Attributes' \subseteq \{ a | (u,a) \in \UA \}$.

%The encoding of a traditional \gls{abac} policies is therefore:
%\begin{center}
%	$(\underline{\user}) \in \Users$ \\
%	$(\underline{\attribute}) \in \Attributes$ \\
%	$(\underline{\resource}) \in \Resources$ \\
%	$(\underline{\user}, \underline{\attribute}, \attributeValue) \in \UA$ \\
%	$(\underline{\resource}, \abeAccessStructure, \underline{\operation} ) \in \FA$
%\end{center}

%See the state-change rules of \gls{abac} in \Cref{table:abacstatechangerules}
%\begin{table*}[!b]
%	\rowcolors{0}{white}{cloud}
%	\caption{\gls{abac} state-change rules and entailment relation}
%	\begin{tabularx}{\linewidth}{*1{>{\sloppy\arraybackslash} X }|c}
%		
%		\toprule
%		
%		\multicolumn{1}{c|}{Operation} &
%		Resulting \acrshort{abac} Policy State \\
%		
%		
%		\hline
%		
%		
%		addUser($\user$) & 
%		$\langle \Users \cup \{\user\}, \Attributes, \Resources, \UA, \FA \rangle$ \\
%		
%		deleteUser($\user$) & 
%		$\langle \Users \setminus \{\user\}, \Attributes, \Resources, \UA, \FA \rangle$ \\
%		
%		addAttribute($\attribute$) & 
%		$\langle \Users, \Attributes \cup \{\attribute\}, \Resources, \UA, \FA \rangle$ \\
%		
%		deleteAttribute($\attribute$) & 
%		$\langle \Users, \Attributes \setminus \{\attribute\}, \Resources, \UA, \FA \rangle$ \\
%		
%		addResource($\resource$) & 
%		$\langle \Users, \Attributes, \Resources \cup \{\resource\}, \UA, \FA \rangle$ \\
%		
%		deleteResource($\resource$) & 
%		$\langle \Users, \Attributes, \Resources \setminus \{\resource\}, \UA, \FA \rangle$ \\
%		
%		assignAttributeToUser($\user, \attribute, \attributeValue$) & 
%		$\langle \Users, \Attributes, \Resources, \UA \cup \{(\user, \attribute, \attributeValue)\}, \FA \rangle$ \\
%		
%		revokeAttributeFromUser($\user, \attribute$) & 
%		$\langle \Users, \Attributes, \Resources, \UA \setminus  \{(\user, \attribute, \wildcard)\}, \FA \rangle$ \\
%		
%		assignAccessStructureToResource($\langle \resource, \operation \rangle, \abeAccessStructure$) & 
%		$\langle \Users, \Attributes, \Resources, \UA, \FA \cup \{(\langle \resource, \operation \rangle, \abeAccessStructure)\} \rangle$ \\
%		
%		revokeAccessStructurefromResource($\langle \resource, \operation \rangle$) & 
%		$\langle \Users, \Attributes, \Resources, \UA, \FA \setminus \{(\langle \resource, \operation \rangle, \wildcard)\} \rangle$ \\
%		
%		
%		\hline
%		
%		
%		\rowcolor{white}
%		\multicolumn{2}{l}{$\gamma \vdash(\langle \user, \langle \resource, \operation \rangle \rangle): \mathit{true} \iff \exists (\langle \resource, \operation \rangle, \abeAccessStructure) \in \FA \mid~\abeAccessStructure \subseteq \{ (\attribute, \attributeValue) | (\user, \attribute, \attributeValue) \in \UA \}$} \\
%		
%		
%		\bottomrule
%		
%	\end{tabularx}
%	\label{table:abacstatechangerules}
%\end{table*}


\vspace{-1.5mm}
\newpage
